% Options for packages loaded elsewhere
\PassOptionsToPackage{unicode}{hyperref}
\PassOptionsToPackage{hyphens}{url}
\PassOptionsToPackage{dvipsnames,svgnames,x11names}{xcolor}
\documentclass[
  spanish,
  10pt,
  a4paper,
]{article}
\usepackage{xcolor}
\usepackage[margin=20mm,headheight=15pt]{geometry}
\usepackage{amsmath,amssymb}
\setcounter{secnumdepth}{-\maxdimen} % remove section numbering
\usepackage{iftex}
\ifPDFTeX
  \usepackage[T1]{fontenc}
  \usepackage[utf8]{inputenc}
  \usepackage{textcomp} % provide euro and other symbols
\else % if luatex or xetex
  \usepackage{unicode-math} % this also loads fontspec
  \defaultfontfeatures{Scale=MatchLowercase}
  \defaultfontfeatures[\rmfamily]{Ligatures=TeX,Scale=1}
\fi
\usepackage{lmodern}
\ifPDFTeX\else
  % xetex/luatex font selection
  \setmainfont[]{DejaVu Sans}
  \setmonofont[]{DejaVu Sans Mono}
\fi
% Use upquote if available, for straight quotes in verbatim environments
\IfFileExists{upquote.sty}{\usepackage{upquote}}{}
\IfFileExists{microtype.sty}{% use microtype if available
  \usepackage[]{microtype}
  \UseMicrotypeSet[protrusion]{basicmath} % disable protrusion for tt fonts
}{}
\makeatletter
\@ifundefined{KOMAClassName}{% if non-KOMA class
  \IfFileExists{parskip.sty}{%
    \usepackage{parskip}
  }{% else
    \setlength{\parindent}{0pt}
    \setlength{\parskip}{6pt plus 2pt minus 1pt}}
}{% if KOMA class
  \KOMAoptions{parskip=half}}
\makeatother
\usepackage{color}
\usepackage{fancyvrb}
\newcommand{\VerbBar}{|}
\newcommand{\VERB}{\Verb[commandchars=\\\{\}]}
\DefineVerbatimEnvironment{Highlighting}{Verbatim}{commandchars=\\\{\}}
% Add ',fontsize=\small' for more characters per line
\newenvironment{Shaded}{}{}
\newcommand{\AlertTok}[1]{\textbf{#1}}
\newcommand{\AnnotationTok}[1]{\textit{#1}}
\newcommand{\AttributeTok}[1]{#1}
\newcommand{\BaseNTok}[1]{#1}
\newcommand{\BuiltInTok}[1]{#1}
\newcommand{\CharTok}[1]{#1}
\newcommand{\CommentTok}[1]{\textit{#1}}
\newcommand{\CommentVarTok}[1]{\textit{#1}}
\newcommand{\ConstantTok}[1]{#1}
\newcommand{\ControlFlowTok}[1]{\textbf{#1}}
\newcommand{\DataTypeTok}[1]{\underline{#1}}
\newcommand{\DecValTok}[1]{#1}
\newcommand{\DocumentationTok}[1]{\textit{#1}}
\newcommand{\ErrorTok}[1]{\textbf{#1}}
\newcommand{\ExtensionTok}[1]{#1}
\newcommand{\FloatTok}[1]{#1}
\newcommand{\FunctionTok}[1]{#1}
\newcommand{\ImportTok}[1]{#1}
\newcommand{\InformationTok}[1]{\textit{#1}}
\newcommand{\KeywordTok}[1]{\textbf{#1}}
\newcommand{\NormalTok}[1]{#1}
\newcommand{\OperatorTok}[1]{#1}
\newcommand{\OtherTok}[1]{#1}
\newcommand{\PreprocessorTok}[1]{\textbf{#1}}
\newcommand{\RegionMarkerTok}[1]{#1}
\newcommand{\SpecialCharTok}[1]{#1}
\newcommand{\SpecialStringTok}[1]{#1}
\newcommand{\StringTok}[1]{#1}
\newcommand{\VariableTok}[1]{#1}
\newcommand{\VerbatimStringTok}[1]{#1}
\newcommand{\WarningTok}[1]{\textit{#1}}
\usepackage{longtable,booktabs,array}
\usepackage{calc} % for calculating minipage widths
% Correct order of tables after \paragraph or \subparagraph
\usepackage{etoolbox}
\makeatletter
\patchcmd\longtable{\par}{\if@noskipsec\mbox{}\fi\par}{}{}
\makeatother
% Allow footnotes in longtable head/foot
\IfFileExists{footnotehyper.sty}{\usepackage{footnotehyper}}{\usepackage{footnote}}
\makesavenoteenv{longtable}
\ifLuaTeX
\usepackage[bidi=basic,shorthands=off,]{babel}
\else
\usepackage[bidi=default,shorthands=off,]{babel}
\fi
\ifPDFTeX
\else
\babelfont{rm}[]{DejaVu Sans}
\fi
\ifLuaTeX
  \usepackage{selnolig} % disable illegal ligatures
\fi
\setlength{\emergencystretch}{3em} % prevent overfull lines
\providecommand{\tightlist}{%
  \setlength{\itemsep}{0pt}\setlength{\parskip}{0pt}}
\usepackage{xcolor}
\definecolor{NavyBI}{HTML}{173C59}
\definecolor{TealBI}{HTML}{087C85}
\definecolor{PaleBI}{HTML}{EFF5F8}
\definecolor{LineBI}{HTML}{D5E3EA}
\usepackage{titlesec}
\titleformat{\section}{\color{NavyBI}\sffamily\Large\bfseries}{\thesection}{.55em}{}
\titleformat{\subsection}{\color{TealBI}\sffamily\normalsize\bfseries}{\thesubsection}{.5em}{}
\titlespacing*{\section}{0pt}{3pt}{10pt}
\titlespacing*{\subsection}{0pt}{10pt}{5pt}
\usepackage[most]{tcolorbox}
\usepackage{fvextra}
\usepackage{framed}
\definecolor{shadecolor}{HTML}{EFF5F8}
\AtBeginDocument{%
\RecustomVerbatimEnvironment{Highlighting}{Verbatim}{commandchars=\\\{\},breaklines=true,breakanywhere=true,breaksymbolleft={},fontsize=\footnotesize}%
\renewenvironment{Shaded}{\begin{snugshade}}{\end{snugshade}}%
}
\usepackage{fancyhdr}
\usepackage{lastpage}
\pagestyle{fancy}
\fancyhf{}
\fancyhead[L]{\color{TealBI}\sffamily\fontsize{8}{10}\selectfont BI PORTABLE · GUÍA DE INSTALACIÓN}
\fancyhead[R]{\color{NavyBI}\sffamily\fontsize{8}{10}\selectfont 25 NOV 2026}
\fancyfoot[L]{\color{NavyBI}\sffamily\fontsize{7.5}{9}\selectfont Nueva entrega pendiente de construcción y aceptación}
\fancyfoot[R]{\color{NavyBI}\sffamily\fontsize{8}{10}\selectfont\thepage\ / \pageref{LastPage}}
\renewcommand{\headrulewidth}{.4pt}
\renewcommand{\footrulewidth}{.2pt}
\setlength{\emergencystretch}{3em}
\setlength{\parskip}{5pt}
\setlength{\parindent}{0pt}
\AtBeginDocument{\hypersetup{pdftitle={Actualizar BI Portable: guía para las programadoras},pdfauthor={Proyecto BI + IA},pdfsubject={Instalación limpia futura; comandos manuales; desarrollo y release}}}
\usepackage{bookmark}
\IfFileExists{xurl.sty}{\usepackage{xurl}}{} % add URL line breaks if available
\urlstyle{same}
\hypersetup{
  pdflang={es-EC},
  colorlinks=true,
  linkcolor={NavyBI},
  filecolor={Maroon},
  citecolor={Blue},
  urlcolor={TealBI},
  pdfcreator={LaTeX via pandoc}}

\author{}
\date{}

\begin{document}

\begin{titlepage}
\vspace*{12mm}
{\color{TealBI}\sffamily\bfseries PROYECTO BI + IA\par}
\vspace{10mm}
{\color{NavyBI}\sffamily\bfseries\fontsize{29}{34}\selectfont Actualizar BI Portable\par}
\vspace{5mm}
{\Large Guía de instalación limpia para las programadoras\par}
\vspace{10mm}
{\large Desarrollo y release · Una orden a la vez\par}
\vspace{12mm}
\begin{tcolorbox}[colback=PaleBI,colframe=TealBI,boxrule=.6pt,arc=2mm]
\textbf{Esta guía prepara la siguiente entrega.} El paquete nuevo todavía debe construirse, recibir su nombre definitivo y superar sus verificaciones. Esta guía no convierte el portable del 17 de noviembre en una versión actualizada.
\end{tcolorbox}
\vspace{8mm}
La instalación se realizará en el computador receptor. Los contenedores y datos del computador del tutor no se retirarán ni se copiarán por este procedimiento.
\vfill
{\color{TealBI}\bfseries Edición documental: 25 de noviembre de 2026\par}
\vspace{4mm}
El PDF explica el recorrido. El archivo Markdown conserva cada comando completo para copiarlo sin los saltos de línea de la página impresa. Ambos son documentación, no instaladores ni scripts.
\end{titlepage}

\tableofcontents
\clearpage

\section{Antes de empezar}\label{antes-de-empezar}

El resultado buscado es una \textbf{instalación limpia} con el código
nuevo, los servicios de desarrollo y release, la estructura PostgreSQL
completa, un administrador y las dos bases públicas de demostración en
SQL Server. Las API keys, conexiones, propuestas, ejecuciones y
datamarts privados del tutor no se trasladan.

Cada ambiente conserva sus propias bases, credenciales y clave de
cifrado. Con \textbf{8 GB de RAM se utilizará un solo ambiente a la
vez}. No iniciar Ollama ni otra instancia SQL adicional durante la
demostración.

\subsection{Cómo usar esta guía}\label{cuxf3mo-usar-esta-guuxeda}

\begin{enumerate}
\def\labelenumi{\arabic{enumi}.}
\tightlist
\item
  Leer y verificar el paquete nuevo antes de retirar el anterior.
\item
  Identificar los recursos antiguos por sus etiquetas Docker.
\item
  Detener y eliminar sólo los recursos BI comprobados.
\item
  Iniciar desarrollo y comprobarlo; detenerlo antes de probar release.
\item
  Configurar cada ambiente y repetir el recorrido BI con ambas fuentes.
\end{enumerate}

\textbf{No copie bloques completos a la terminal.} Ejecute una orden,
compruebe el resultado y continúe. Los comandos que comienzan con
\texttt{read} piden escribir un nombre real observado. No escriba el
símbolo \texttt{\$} del indicador de la terminal.

\subsection{Comprobar el Ubuntu
receptor}\label{comprobar-el-ubuntu-receptor}

Si trabaja en Windows, lea el siguiente apartado de preparación WSL y
continúe con la sección común del paquete nuevo. La identidad
\texttt{ls-va@Lenovo} corresponde únicamente al Ubuntu de la
programadora que reemplazará su instalación anterior.

\begin{Shaded}
\begin{Highlighting}[]
\FunctionTok{whoami}
\end{Highlighting}
\end{Shaded}

\begin{Shaded}
\begin{Highlighting}[]
\BuiltInTok{printf} \StringTok{\textquotesingle{}\%s\textbackslash{}n\textquotesingle{}} \StringTok{"}\VariableTok{$HOME}\StringTok{"}
\end{Highlighting}
\end{Shaded}

\begin{Shaded}
\begin{Highlighting}[]
\FunctionTok{hostname} \AttributeTok{{-}s}
\end{Highlighting}
\end{Shaded}

En el Ubuntu conocido deben aparecer \textbf{ls-va},
\textbf{/home/ls-va} y \textbf{Lenovo}. Si aparece \texttt{ceo},
\texttt{/home/ceo} o un equipo distinto, detenerse. Las órdenes de
retirada incluyen esa comprobación para que no operen en el equipo del
tutor. No cambie la comprobación para hacer desaparecer un error.

\begin{Shaded}
\begin{Highlighting}[]
\FunctionTok{uname} \AttributeTok{{-}m}
\end{Highlighting}
\end{Shaded}

\begin{Shaded}
\begin{Highlighting}[]
\FunctionTok{free} \AttributeTok{{-}h}
\end{Highlighting}
\end{Shaded}

\begin{Shaded}
\begin{Highlighting}[]
\FunctionTok{df} \AttributeTok{{-}h} \StringTok{"}\VariableTok{$HOME}\StringTok{"}
\end{Highlighting}
\end{Shaded}

\begin{Shaded}
\begin{Highlighting}[]
\ExtensionTok{docker}\NormalTok{ version}
\end{Highlighting}
\end{Shaded}

\begin{Shaded}
\begin{Highlighting}[]
\ExtensionTok{docker}\NormalTok{ compose version}
\end{Highlighting}
\end{Shaded}

Se requiere arquitectura \texttt{x86\_64}, Docker con cliente y servidor
disponibles y el complemento Compose. Hay que disponer de espacio para
el archivo recibido, las imágenes importadas y los volúmenes nuevos.
Conservar Docker ya instalado; un fallo de la aplicación no exige
reinstalarlo.

\clearpage

\section{Preparar Windows para una instalación
inicial}\label{preparar-windows-para-una-instalaciuxf3n-inicial}

Este recorrido corresponde al equipo Windows de la otra programadora.
\textbf{No se ha probado su instalación remota.} Se prepara Ubuntu bajo
WSL 2 y después se utilizan los mismos archivos portable desde su
carpeta Linux. Si Docker y WSL ya funcionan, comprobarlos y
conservarlos; no reinstalarlos por rutina.

\subsection{Comprobar Windows y WSL desde
PowerShell}\label{comprobar-windows-y-wsl-desde-powershell}

Abrir PowerShell y ejecutar una orden a la vez:

\begin{Shaded}
\begin{Highlighting}[]
\NormalTok{winver}
\end{Highlighting}
\end{Shaded}

\begin{Shaded}
\begin{Highlighting}[]
\NormalTok{wsl }\OperatorTok{{-}{-}}\NormalTok{version}
\end{Highlighting}
\end{Shaded}

\begin{Shaded}
\begin{Highlighting}[]
\NormalTok{wsl }\OperatorTok{{-}}\NormalTok{l }\OperatorTok{{-}}\NormalTok{v}
\end{Highlighting}
\end{Shaded}

La documentación Docker consultada el 25-11-2026 enumera Windows 10 Pro
de 64 bits, versión 22H2, compilación 19045, WSL 2.1.5 o posterior y 8
GB de RAM para esta ruta. Confirmar además las condiciones vigentes de
soporte del sistema operativo en la página oficial antes de elegir el
instalador.

\subsection{Instalar WSL sólo si
falta}\label{instalar-wsl-suxf3lo-si-falta}

Abrir \textbf{PowerShell como administrador}. Si Ubuntu 24.04 no está
instalado, utilizar:

\begin{Shaded}
\begin{Highlighting}[]
\NormalTok{wsl }\OperatorTok{{-}{-}}\NormalTok{install }\OperatorTok{{-}}\NormalTok{d Ubuntu{-}24}\OperatorTok{.}\DecValTok{04}
\end{Highlighting}
\end{Shaded}

Reiniciar Windows si se solicita. Abrir \textbf{Ubuntu 24.04} desde
Inicio, esperar su preparación y crear el usuario y contraseña Linux. La
contraseña no se muestra mientras se escribe. No es la contraseña de la
aplicación BI.

Si WSL ya existe y requiere actualización, utilizar en PowerShell:

\begin{Shaded}
\begin{Highlighting}[]
\NormalTok{wsl }\OperatorTok{{-}{-}}\NormalTok{update}
\end{Highlighting}
\end{Shaded}

\begin{Shaded}
\begin{Highlighting}[]
\NormalTok{wsl }\OperatorTok{{-}}\NormalTok{l }\OperatorTok{{-}}\NormalTok{v}
\end{Highlighting}
\end{Shaded}

La distribución elegida debe indicar \textbf{VERSION 2}. Si esa
distribución exacta aparece como \texttt{Ubuntu-24.04} con versión 1,
convertirla con:

\begin{Shaded}
\begin{Highlighting}[]
\NormalTok{wsl }\OperatorTok{{-}{-}}\NormalTok{set{-}version Ubuntu{-}24}\OperatorTok{.}\DecValTok{04} \DecValTok{2}
\end{Highlighting}
\end{Shaded}

No cambiar el nombre de una distribución distinta por aproximación. Si
falla la virtualización, resolver ese requisito con el responsable del
equipo antes de continuar.

\clearpage

\subsection{Configurar Docker Desktop y recibir el
archivo}\label{configurar-docker-desktop-y-recibir-el-archivo}

Descargar Docker Desktop para Windows \textbf{x86\_64} desde la página
oficial Docker, no desde sitios de terceros. Instalarlo sólo si falta y
abrirlo; completar el reinicio solicitado, si corresponde.

En Docker Desktop:

\begin{enumerate}
\def\labelenumi{\arabic{enumi}.}
\tightlist
\item
  En \textbf{Settings → General}, utilizar el motor basado en WSL 2.
\item
  En \textbf{Settings → Resources → WSL Integration}, habilitar la
  distribución Ubuntu utilizada y aplicar los cambios.
\item
  Utilizar \textbf{contenedores Linux}. Si el menú ofrece cambiar a
  Linux containers, elegir esa opción.
\item
  Esperar a que Docker indique que su motor está disponible.
\end{enumerate}

Las órdenes siguientes se escriben en la terminal \textbf{Ubuntu}, no en
PowerShell:

\begin{Shaded}
\begin{Highlighting}[]
\ExtensionTok{docker}\NormalTok{ version}
\end{Highlighting}
\end{Shaded}

\begin{Shaded}
\begin{Highlighting}[]
\ExtensionTok{docker}\NormalTok{ compose version}
\end{Highlighting}
\end{Shaded}

\begin{Shaded}
\begin{Highlighting}[]
\ExtensionTok{docker}\NormalTok{ info }\AttributeTok{{-}{-}format} \StringTok{\textquotesingle{}\{\{.OSType\}\}\textquotesingle{}}
\end{Highlighting}
\end{Shaded}

Deben responder el cliente y servidor, Compose y el tipo \texttt{linux}.
No instalar otro servidor Docker dentro de Ubuntu si se está utilizando
correctamente la integración de Docker Desktop.

Preparar una carpeta de recepción en el HOME Linux:

\begin{Shaded}
\begin{Highlighting}[]
\FunctionTok{mkdir} \AttributeTok{{-}p} \StringTok{"}\VariableTok{$HOME}\StringTok{/Descargas"}
\end{Highlighting}
\end{Shaded}

\begin{Shaded}
\begin{Highlighting}[]
\BuiltInTok{cd} \StringTok{"}\VariableTok{$HOME}\StringTok{/Descargas"}
\end{Highlighting}
\end{Shaded}

\begin{Shaded}
\begin{Highlighting}[]
\ExtensionTok{explorer.exe}\NormalTok{ .}
\end{Highlighting}
\end{Shaded}

Se abrirá el Explorador de Windows en esa carpeta Ubuntu. Copiar allí el
\textbf{nuevo} \texttt{.tar.gz} y su \texttt{.tar.gz.sha256} desde la
ubicación donde se recibieron. No extraer el desarrollo directamente en
una ruta \texttt{C:\textbackslash{}...}; la instalación se mantendrá en
\texttt{\textasciitilde{}/BI/} dentro de Ubuntu.

Continuar con \textbf{Preparar el paquete nuevo sin arrancarlo} y
\textbf{Identificar nombres, imágenes y proyectos}. Al ser una
instalación inicial, \textbf{saltar los dos apartados de retirada de
contenedores, volúmenes y carpetas Ubuntu Lenovo}. Seguir con Iniciar
desarrollo. Con 8 GB, probar un solo ambiente a la vez.

Fuentes oficiales, consulta del 25-11-2026:
\href{https://docs.docker.com/desktop/setup/install/windows-install/}{Docker
Desktop para Windows} y
\href{https://learn.microsoft.com/en-us/windows/wsl/install}{Microsoft:
instalar WSL}.

\clearpage

\section{Preparar el paquete nuevo sin
arrancarlo}\label{preparar-el-paquete-nuevo-sin-arrancarlo}

\textbf{Este paso sólo se realizará cuando el tutor entregue el archivo
nuevo y su suma SHA-256.} No seleccionar el archivo del 17 de noviembre
suponiendo que contiene estas correcciones. El manifiesto nuevo deberá
identificar el corte real del código y las imágenes de ambos modos.

Copiar juntos el \texttt{.tar.gz} y el \texttt{.tar.gz.sha256} a
Descargas. Abrir una terminal y mantenerla durante la instalación.

\begin{Shaded}
\begin{Highlighting}[]
\BuiltInTok{cd} \StringTok{"}\VariableTok{$HOME}\StringTok{/Descargas"}
\end{Highlighting}
\end{Shaded}

\begin{Shaded}
\begin{Highlighting}[]
\FunctionTok{ls} \AttributeTok{{-}lh}\NormalTok{ BI\_portable\_}\PreprocessorTok{*}\NormalTok{.tar.gz BI\_portable\_}\PreprocessorTok{*}\NormalTok{.tar.gz.sha256}
\end{Highlighting}
\end{Shaded}

Escribir el nombre exacto del paquete nuevo, \textbf{sin} la terminación
\texttt{.tar.gz}. Esta variable evita inventar una fecha o etiqueta
futura.

\begin{Shaded}
\begin{Highlighting}[]
\BuiltInTok{read} \AttributeTok{{-}r} \AttributeTok{{-}p} \StringTok{"Nombre real del paquete nuevo, sin .tar.gz: "} \VariableTok{BI\_PACKAGE}
\end{Highlighting}
\end{Shaded}

\begin{Shaded}
\begin{Highlighting}[]
\FunctionTok{sha256sum} \AttributeTok{{-}c} \StringTok{"}\VariableTok{$BI\_PACKAGE}\StringTok{.tar.gz.sha256"}
\end{Highlighting}
\end{Shaded}

Debe indicar \texttt{OK}. Si no coincide, copiar nuevamente los archivos
y detener aquí la instalación.

\begin{Shaded}
\begin{Highlighting}[]
\FunctionTok{mkdir} \AttributeTok{{-}p} \StringTok{"}\VariableTok{$HOME}\StringTok{/BI"}
\end{Highlighting}
\end{Shaded}

\begin{Shaded}
\begin{Highlighting}[]
\BuiltInTok{test}\NormalTok{ ! }\AttributeTok{{-}e} \StringTok{"}\VariableTok{$HOME}\StringTok{/BI/}\VariableTok{$BI\_PACKAGE}\StringTok{"} \KeywordTok{\&\&} \FunctionTok{tar} \AttributeTok{{-}{-}keep{-}old{-}files} \AttributeTok{{-}xzf} \StringTok{"}\VariableTok{$HOME}\StringTok{/Descargas/}\VariableTok{$BI\_PACKAGE}\StringTok{.tar.gz"} \AttributeTok{{-}C} \StringTok{"}\VariableTok{$HOME}\StringTok{/BI"}
\end{Highlighting}
\end{Shaded}

Si la carpeta existe, no extraer encima. Una carpeta usada por
desarrollo no debe reemplazarse mientras sus servicios están activos.

\begin{Shaded}
\begin{Highlighting}[]
\VariableTok{BI\_NEW\_DIRECTORY}\OperatorTok{=}\StringTok{"}\VariableTok{$HOME}\StringTok{/BI/}\VariableTok{$BI\_PACKAGE}\StringTok{"}
\end{Highlighting}
\end{Shaded}

\begin{Shaded}
\begin{Highlighting}[]
\BuiltInTok{cd} \StringTok{"}\VariableTok{$BI\_NEW\_DIRECTORY}\StringTok{"}
\end{Highlighting}
\end{Shaded}

\begin{Shaded}
\begin{Highlighting}[]
\BuiltInTok{pwd} \AttributeTok{{-}P}
\end{Highlighting}
\end{Shaded}

\begin{Shaded}
\begin{Highlighting}[]
\FunctionTok{sha256sum} \AttributeTok{{-}c}\NormalTok{ SHA256SUMS.txt}
\end{Highlighting}
\end{Shaded}

Todas las sumas deben coincidir. Todavía no editar el código ni las
configuraciones.

\begin{Shaded}
\begin{Highlighting}[]
\FunctionTok{chmod}\NormalTok{ 600 config/desarrollo.env config/release.env CREDENCIALES\_LOCALES.txt}
\end{Highlighting}
\end{Shaded}

\begin{Shaded}
\begin{Highlighting}[]
\FunctionTok{less}\NormalTok{ MANIFIESTO.json}
\end{Highlighting}
\end{Shaded}

Salir de \texttt{less} con la tecla \textbf{q}. Comprobar que el
manifiesto corresponde al paquete nuevo, explica si es una entrega
limpia y registra las pruebas realizadas. Los resultados de noviembre
anterior no sustituyen sus verificaciones nuevas.

\clearpage

\section{Identificar nombres, imágenes y
proyectos}\label{identificar-nombres-imuxe1genes-y-proyectos}

Desde la carpeta nueva, validar ambos Compose sin imprimir sus variables
privadas.

\begin{Shaded}
\begin{Highlighting}[]
\ExtensionTok{docker}\NormalTok{ compose }\AttributeTok{{-}{-}env{-}file}\NormalTok{ config/desarrollo.env }\AttributeTok{{-}f}\NormalTok{ compose.desarrollo.yaml config }\AttributeTok{{-}{-}quiet}
\end{Highlighting}
\end{Shaded}

\begin{Shaded}
\begin{Highlighting}[]
\ExtensionTok{docker}\NormalTok{ compose }\AttributeTok{{-}{-}env{-}file}\NormalTok{ config/release.env }\AttributeTok{{-}f}\NormalTok{ compose.release.yaml config }\AttributeTok{{-}{-}quiet}
\end{Highlighting}
\end{Shaded}

\begin{Shaded}
\begin{Highlighting}[]
\FunctionTok{sed} \AttributeTok{{-}n} \StringTok{\textquotesingle{}/\^{}name:/p\textquotesingle{}}\NormalTok{ compose.desarrollo.yaml compose.release.yaml}
\end{Highlighting}
\end{Shaded}

Estos son los nombres de proyecto \textbf{nuevos}. Anotarlos; ninguno se
retirará. No añadir \texttt{-p} con un nombre de una entrega antigua.
Usar el nombre declarado en los archivos recibidos.

\begin{Shaded}
\begin{Highlighting}[]
\ExtensionTok{docker}\NormalTok{ compose }\AttributeTok{{-}{-}env{-}file}\NormalTok{ config/desarrollo.env }\AttributeTok{{-}f}\NormalTok{ compose.desarrollo.yaml config }\AttributeTok{{-}{-}images}
\end{Highlighting}
\end{Shaded}

\begin{Shaded}
\begin{Highlighting}[]
\ExtensionTok{docker}\NormalTok{ compose }\AttributeTok{{-}{-}env{-}file}\NormalTok{ config/release.env }\AttributeTok{{-}f}\NormalTok{ compose.release.yaml config }\AttributeTok{{-}{-}images}
\end{Highlighting}
\end{Shaded}

\begin{Shaded}
\begin{Highlighting}[]
\ExtensionTok{docker}\NormalTok{ load }\AttributeTok{{-}i}\NormalTok{ imagenes.tar}
\end{Highlighting}
\end{Shaded}

El archivo debe cargar las imágenes exactas del manifiesto: backend y
frontend para desarrollo y release, PostgreSQL y SQL Server. No usar
\texttt{-\/-build} ni descargar imágenes para el primer inicio.

\subsection{Descubrir el BI anterior aunque su carpeta haya
desaparecido}\label{descubrir-el-bi-anterior-aunque-su-carpeta-haya-desaparecido}

\begin{Shaded}
\begin{Highlighting}[]
\ExtensionTok{docker}\NormalTok{ compose ls }\AttributeTok{{-}{-}all}
\end{Highlighting}
\end{Shaded}

\begin{Shaded}
\begin{Highlighting}[]
\ExtensionTok{docker}\NormalTok{ ps }\AttributeTok{{-}a} \AttributeTok{{-}{-}format} \StringTok{\textquotesingle{}table \{\{.Names\}\}\textbackslash{}t\{\{.Image\}\}\textbackslash{}t\{\{.Status\}\}\textbackslash{}t\{\{.Ports\}\}\textquotesingle{}}
\end{Highlighting}
\end{Shaded}

\begin{longtable}[]{@{}
  >{\raggedright\arraybackslash}p{(\linewidth - 2\tabcolsep) * \real{0.5000}}
  >{\raggedright\arraybackslash}p{(\linewidth - 2\tabcolsep) * \real{0.5000}}@{}}
\toprule\noalign{}
\begin{minipage}[b]{\linewidth}\raggedright
Instalación anterior conocida
\end{minipage} & \begin{minipage}[b]{\linewidth}\raggedright
Nombre de proyecto que puede aparecer
\end{minipage} \\
\midrule\noalign{}
\endhead
\bottomrule\noalign{}
\endlastfoot
Portable de octubre & \texttt{bi-limpio-dev} y
\texttt{bi-limpio-release} \\
Portable del 17 de noviembre & \texttt{bi-portable-20261117-dev} y
\texttt{bi-portable-20261117-release} \\
Clones antiguos & \texttt{bi-ia-prototype} y
\texttt{bi-ia-prototype-release} \\
\end{longtable}

Eliminar una carpeta con Supr \textbf{no detiene ni elimina sus
contenedores o volúmenes}. Por eso puede quedar ocupado el puerto 11433
aunque el directorio ya no exista. Se puede identificar y retirar el
contenedor por sus etiquetas; no hace falta reconstruir el YAML antiguo
ni utilizar el nuevo para borrar el anterior.

\texttt{bd-stack}, MySQL, MariaDB, pgAdmin, otros proyectos y recursos
sin pertenencia BI comprobada quedan fuera de la sustitución.

\clearpage

\section{Retirar contenedores antiguos, uno por
uno}\label{retirar-contenedores-antiguos-uno-por-uno}

Este paso es exclusivo del Ubuntu receptor identificado. La sustitución
limpia del BI antiguo fue autorizada sin respaldo. Sus configuraciones y
datos de prueba se perderán cuando se retiren sus volúmenes. \textbf{No
se ejecutaron estas retiradas durante la elaboración de la guía.}

Elegir un proyecto antiguo de la tabla anterior. Repetir después el
procedimiento para cada proyecto que realmente exista, incluido release.

\begin{Shaded}
\begin{Highlighting}[]
\BuiltInTok{read} \AttributeTok{{-}r} \AttributeTok{{-}p} \StringTok{"Proyecto BI antiguo exacto: "} \VariableTok{BI\_OLD\_PROJECT}
\end{Highlighting}
\end{Shaded}

\begin{Shaded}
\begin{Highlighting}[]
\ExtensionTok{docker}\NormalTok{ ps }\AttributeTok{{-}a} \AttributeTok{{-}{-}filter} \StringTok{"label=com.docker.compose.project=}\VariableTok{$BI\_OLD\_PROJECT}\StringTok{"} \AttributeTok{{-}{-}format} \StringTok{\textquotesingle{}table \{\{.Names\}\}\textbackslash{}t\{\{.Status\}\}\textbackslash{}t\{\{.Ports\}\}\textquotesingle{}}
\end{Highlighting}
\end{Shaded}

\begin{Shaded}
\begin{Highlighting}[]
\ExtensionTok{docker}\NormalTok{ volume ls }\AttributeTok{{-}{-}filter} \StringTok{"label=com.docker.compose.project=}\VariableTok{$BI\_OLD\_PROJECT}\StringTok{"}
\end{Highlighting}
\end{Shaded}

Elegir \textbf{un} nombre de contenedor del listado. Por ejemplo, sólo
si aparece: \texttt{bi-limpio-dev-sqlserver-1}. No escribir un nombre de
la instalación nueva.

\begin{Shaded}
\begin{Highlighting}[]
\BuiltInTok{read} \AttributeTok{{-}r} \AttributeTok{{-}p} \StringTok{"Nombre exacto de UN contenedor BI antiguo listado: "} \VariableTok{BI\_OLD\_CONTAINER}
\end{Highlighting}
\end{Shaded}

\begin{Shaded}
\begin{Highlighting}[]
\ExtensionTok{docker}\NormalTok{ container inspect }\StringTok{"}\VariableTok{$BI\_OLD\_CONTAINER}\StringTok{"} \AttributeTok{{-}{-}format} \StringTok{\textquotesingle{}\{\{.Name\}\} | proyecto=\{\{index .Config.Labels "com.docker.compose.project"\}\} | carpeta=\{\{index .Config.Labels "com.docker.compose.project.working\_dir"\}\}\textquotesingle{}}
\end{Highlighting}
\end{Shaded}

Comprobar proyecto y carpeta. Las rutas antiguas conocidas son:

\begin{itemize}
\tightlist
\item
  \texttt{\textasciitilde{}/Descargas/BI\_portable\_2026-10-29\_v02}.
\item
  \texttt{\textasciitilde{}/BI/BI\_portable\_2026-11-17\_v01}.
\item
  Los dos clones BI de \texttt{\textasciitilde{}/Proyectos/}.
\end{itemize}

Que la carpeta ya no exista no cambia la etiqueta histórica del
contenedor.

La orden siguiente sólo admite los proyectos antiguos de la tabla, exige
el equipo receptor y vuelve a comprobar la etiqueta del contenedor antes
de detenerlo:

\begin{Shaded}
\begin{Highlighting}[]
\BuiltInTok{test} \StringTok{"}\VariableTok{$(}\FunctionTok{id} \AttributeTok{{-}un}\VariableTok{)}\StringTok{:}\VariableTok{$HOME}\StringTok{:}\VariableTok{$(}\FunctionTok{hostname} \AttributeTok{{-}s}\VariableTok{)}\StringTok{"}\NormalTok{ = }\StringTok{\textquotesingle{}ls{-}va:/home/ls{-}va:Lenovo\textquotesingle{}} \KeywordTok{\&\&} \KeywordTok{[[} \StringTok{"}\VariableTok{$BI\_OLD\_PROJECT}\StringTok{"} \OtherTok{=\textasciitilde{}} \PreprocessorTok{\^{}}\OperatorTok{(}\SpecialStringTok{bi{-}limpio{-}}\OperatorTok{(}\SpecialStringTok{dev}\PreprocessorTok{|}\SpecialStringTok{release}\OperatorTok{)}\PreprocessorTok{|}\SpecialStringTok{bi{-}portable{-}20261117{-}}\OperatorTok{(}\SpecialStringTok{dev}\PreprocessorTok{|}\SpecialStringTok{release}\OperatorTok{)}\PreprocessorTok{|}\SpecialStringTok{bi{-}ia{-}prototype}\OperatorTok{(}\SpecialStringTok{{-}release}\OperatorTok{)}\PreprocessorTok{?}\OperatorTok{)$} \KeywordTok{]]} \KeywordTok{\&\&} \BuiltInTok{test} \StringTok{"}\VariableTok{$(}\ExtensionTok{docker}\NormalTok{ container inspect }\StringTok{"}\VariableTok{$BI\_OLD\_CONTAINER}\StringTok{"} \AttributeTok{{-}{-}format} \StringTok{\textquotesingle{}\{\{index .Config.Labels "com.docker.compose.project"\}\}\textquotesingle{}}\VariableTok{)}\StringTok{"}\NormalTok{ = }\StringTok{"}\VariableTok{$BI\_OLD\_PROJECT}\StringTok{"} \KeywordTok{\&\&} \ExtensionTok{docker}\NormalTok{ stop }\StringTok{"}\VariableTok{$BI\_OLD\_CONTAINER}\StringTok{"}
\end{Highlighting}
\end{Shaded}

Eliminar únicamente ese contenedor, sin forzar:

\begin{Shaded}
\begin{Highlighting}[]
\BuiltInTok{test} \StringTok{"}\VariableTok{$(}\FunctionTok{id} \AttributeTok{{-}un}\VariableTok{)}\StringTok{:}\VariableTok{$HOME}\StringTok{:}\VariableTok{$(}\FunctionTok{hostname} \AttributeTok{{-}s}\VariableTok{)}\StringTok{"}\NormalTok{ = }\StringTok{\textquotesingle{}ls{-}va:/home/ls{-}va:Lenovo\textquotesingle{}} \KeywordTok{\&\&} \KeywordTok{[[} \StringTok{"}\VariableTok{$BI\_OLD\_PROJECT}\StringTok{"} \OtherTok{=\textasciitilde{}} \PreprocessorTok{\^{}}\OperatorTok{(}\SpecialStringTok{bi{-}limpio{-}}\OperatorTok{(}\SpecialStringTok{dev}\PreprocessorTok{|}\SpecialStringTok{release}\OperatorTok{)}\PreprocessorTok{|}\SpecialStringTok{bi{-}portable{-}20261117{-}}\OperatorTok{(}\SpecialStringTok{dev}\PreprocessorTok{|}\SpecialStringTok{release}\OperatorTok{)}\PreprocessorTok{|}\SpecialStringTok{bi{-}ia{-}prototype}\OperatorTok{(}\SpecialStringTok{{-}release}\OperatorTok{)}\PreprocessorTok{?}\OperatorTok{)$} \KeywordTok{]]} \KeywordTok{\&\&} \BuiltInTok{test} \StringTok{"}\VariableTok{$(}\ExtensionTok{docker}\NormalTok{ container inspect }\StringTok{"}\VariableTok{$BI\_OLD\_CONTAINER}\StringTok{"} \AttributeTok{{-}{-}format} \StringTok{\textquotesingle{}\{\{index .Config.Labels "com.docker.compose.project"\}\}\textquotesingle{}}\VariableTok{)}\StringTok{"}\NormalTok{ = }\StringTok{"}\VariableTok{$BI\_OLD\_PROJECT}\StringTok{"} \KeywordTok{\&\&} \ExtensionTok{docker}\NormalTok{ rm }\StringTok{"}\VariableTok{$BI\_OLD\_CONTAINER}\StringTok{"}
\end{Highlighting}
\end{Shaded}

Repetir \textbf{leer nombre → inspeccionar → detener → eliminar} para
los demás contenedores listados de ese proyecto. \texttt{migrate} y
\texttt{seed} pueden estar ya terminados; también deben retirarse si
pertenecen al BI antiguo. Una lista vacía significa que no hay
contenedores de ese proyecto que eliminar.

\clearpage

\section{Retirar volúmenes y carpetas
anteriores}\label{retirar-voluxfamenes-y-carpetas-anteriores}

Conservar el valor de \texttt{BI\_OLD\_PROJECT} comprobado. Si se cerró
la terminal, repetir la identificación anterior; no suponer que las
variables siguen definidas.

\begin{Shaded}
\begin{Highlighting}[]
\ExtensionTok{docker}\NormalTok{ ps }\AttributeTok{{-}a} \AttributeTok{{-}{-}filter} \StringTok{"label=com.docker.compose.project=}\VariableTok{$BI\_OLD\_PROJECT}\StringTok{"}
\end{Highlighting}
\end{Shaded}

El listado de contenedores debe quedar vacío antes de retirar sus
volúmenes.

\begin{Shaded}
\begin{Highlighting}[]
\ExtensionTok{docker}\NormalTok{ volume ls }\AttributeTok{{-}{-}filter} \StringTok{"label=com.docker.compose.project=}\VariableTok{$BI\_OLD\_PROJECT}\StringTok{"}
\end{Highlighting}
\end{Shaded}

\begin{Shaded}
\begin{Highlighting}[]
\BuiltInTok{read} \AttributeTok{{-}r} \AttributeTok{{-}p} \StringTok{"Nombre exacto de UN volumen BI antiguo listado: "} \VariableTok{BI\_OLD\_VOLUME}
\end{Highlighting}
\end{Shaded}

\begin{Shaded}
\begin{Highlighting}[]
\ExtensionTok{docker}\NormalTok{ volume inspect }\StringTok{"}\VariableTok{$BI\_OLD\_VOLUME}\StringTok{"} \AttributeTok{{-}{-}format} \StringTok{\textquotesingle{}\{\{.Name\}\} | proyecto=\{\{index .Labels "com.docker.compose.project"\}\} | volumen=\{\{index .Labels "com.docker.compose.volume"\}\}\textquotesingle{}}
\end{Highlighting}
\end{Shaded}

Sólo si la etiqueta coincide con el proyecto antiguo elegido:

\begin{Shaded}
\begin{Highlighting}[]
\BuiltInTok{test} \StringTok{"}\VariableTok{$(}\FunctionTok{id} \AttributeTok{{-}un}\VariableTok{)}\StringTok{:}\VariableTok{$HOME}\StringTok{:}\VariableTok{$(}\FunctionTok{hostname} \AttributeTok{{-}s}\VariableTok{)}\StringTok{"}\NormalTok{ = }\StringTok{\textquotesingle{}ls{-}va:/home/ls{-}va:Lenovo\textquotesingle{}} \KeywordTok{\&\&} \KeywordTok{[[} \StringTok{"}\VariableTok{$BI\_OLD\_PROJECT}\StringTok{"} \OtherTok{=\textasciitilde{}} \PreprocessorTok{\^{}}\OperatorTok{(}\SpecialStringTok{bi{-}limpio{-}}\OperatorTok{(}\SpecialStringTok{dev}\PreprocessorTok{|}\SpecialStringTok{release}\OperatorTok{)}\PreprocessorTok{|}\SpecialStringTok{bi{-}portable{-}20261117{-}}\OperatorTok{(}\SpecialStringTok{dev}\PreprocessorTok{|}\SpecialStringTok{release}\OperatorTok{)}\PreprocessorTok{|}\SpecialStringTok{bi{-}ia{-}prototype}\OperatorTok{(}\SpecialStringTok{{-}release}\OperatorTok{)}\PreprocessorTok{?}\OperatorTok{)$} \KeywordTok{]]} \KeywordTok{\&\&} \BuiltInTok{test} \StringTok{"}\VariableTok{$(}\ExtensionTok{docker}\NormalTok{ volume inspect }\StringTok{"}\VariableTok{$BI\_OLD\_VOLUME}\StringTok{"} \AttributeTok{{-}{-}format} \StringTok{\textquotesingle{}\{\{index .Labels "com.docker.compose.project"\}\}\textquotesingle{}}\VariableTok{)}\StringTok{"}\NormalTok{ = }\StringTok{"}\VariableTok{$BI\_OLD\_PROJECT}\StringTok{"} \KeywordTok{\&\&} \ExtensionTok{docker}\NormalTok{ volume rm }\StringTok{"}\VariableTok{$BI\_OLD\_VOLUME}\StringTok{"}
\end{Highlighting}
\end{Shaded}

Repetir para cada volumen BI antiguo comprobado, incluido release. Si
Docker informa que está en uso, localizar primero su contenedor; no
añadir una opción de fuerza. Un volumen sin la etiqueta esperada no se
borra por parecido de nombre.

\subsection{Borrar la carpeta sólo después de validar la instalación
nueva}\label{borrar-la-carpeta-suxf3lo-despuuxe9s-de-validar-la-instalaciuxf3n-nueva}

Después de completar los pasos de inicio y acceso de esta guía, abrir el
gestor de archivos y seleccionar \textbf{únicamente} el directorio
antiguo identificado. Revisar su ruta completa y enviarlo a la papelera.
Si ya fue eliminado con Supr, omitir esta operación: la retirada Docker
anterior sigue siendo necesaria.

No borrar \texttt{Descargas}, \texttt{BI}, \texttt{Proyectos} ni
\texttt{/home/ls-va} completos. No borrar la carpeta nueva: desarrollo
utilizará \texttt{fuente/} desde allí. Los clones antiguos
\texttt{\textasciitilde{}/Proyectos/prototipo-bi-ia-dev} y
\texttt{\textasciitilde{}/Proyectos/prototipo-bi-ia-release} se retiran
sólo después de haber comprobado y retirado sus propios recursos Docker.

Las imágenes antiguas que permanezcan son caché y no se ejecutan por sí
solas. No son motivo para realizar una limpieza global; una imagen de
infraestructura puede ser compartida entre etiquetas de entregas. La
guía no utiliza \texttt{docker\ system\ prune},
\texttt{docker\ volume\ prune}, borrados generales ni eliminación
forzada de imágenes.

\clearpage

\section{Iniciar desarrollo y comprobar la
preparación}\label{iniciar-desarrollo-y-comprobar-la-preparaciuxf3n}

Volver a la carpeta \textbf{nueva}. Mantener detenidos los proyectos
viejos y el ambiente release. Los puertos siguientes corresponden al
contrato portable anterior que deberá verificar el manifiesto nuevo.

\begin{Shaded}
\begin{Highlighting}[]
\BuiltInTok{cd} \StringTok{"}\VariableTok{$BI\_NEW\_DIRECTORY}\StringTok{"}
\end{Highlighting}
\end{Shaded}

\begin{Shaded}
\begin{Highlighting}[]
\ExtensionTok{docker}\NormalTok{ compose }\AttributeTok{{-}{-}env{-}file}\NormalTok{ config/desarrollo.env }\AttributeTok{{-}f}\NormalTok{ compose.desarrollo.yaml config }\AttributeTok{{-}{-}quiet}
\end{Highlighting}
\end{Shaded}

\begin{Shaded}
\begin{Highlighting}[]
\ExtensionTok{docker}\NormalTok{ compose }\AttributeTok{{-}{-}env{-}file}\NormalTok{ config/desarrollo.env }\AttributeTok{{-}f}\NormalTok{ compose.desarrollo.yaml up }\AttributeTok{{-}d} \AttributeTok{{-}{-}pull}\NormalTok{ never }\AttributeTok{{-}{-}wait} \AttributeTok{{-}{-}wait{-}timeout}\NormalTok{ 900}
\end{Highlighting}
\end{Shaded}

Esperar a que termine. El primer inicio crea volúmenes vacíos, inicia
PostgreSQL, aplica migraciones y carga la semilla. SQL Server restaura
las dos muestras incluidas en su imagen y configura el lector. No se
importan los volúmenes del tutor.

\begin{Shaded}
\begin{Highlighting}[]
\ExtensionTok{docker}\NormalTok{ compose }\AttributeTok{{-}{-}env{-}file}\NormalTok{ config/desarrollo.env }\AttributeTok{{-}f}\NormalTok{ compose.desarrollo.yaml ps }\AttributeTok{{-}a}
\end{Highlighting}
\end{Shaded}

\begin{longtable}[]{@{}ll@{}}
\toprule\noalign{}
Servicio & Estado esperado \\
\midrule\noalign{}
\endhead
\bottomrule\noalign{}
\endlastfoot
\texttt{postgres}, \texttt{sqlserver}, \texttt{backend},
\texttt{frontend} & Activos y \texttt{healthy} \\
\texttt{migrate}, \texttt{seed} & \texttt{Exited\ (0)}: tareas
terminadas correctamente \\
\end{longtable}

\textbf{Exited (0) no es un error.} \texttt{Exited\ (1)} en
\texttt{migrate} o \texttt{seed} sí exige revisar sus registros antes de
continuar.

\begin{Shaded}
\begin{Highlighting}[]
\ExtensionTok{docker}\NormalTok{ compose }\AttributeTok{{-}{-}env{-}file}\NormalTok{ config/desarrollo.env }\AttributeTok{{-}f}\NormalTok{ compose.desarrollo.yaml exec }\AttributeTok{{-}T}\NormalTok{ backend alembic current}
\end{Highlighting}
\end{Shaded}

El corte actual incorpora la migración \textbf{20261125\_20} y el
parámetro de sesión. El paquete definitivo deberá incluirla o una
revisión posterior declarada. No cambiar manualmente el número en
PostgreSQL ni marcar una migración como aplicada para ocultar un fallo.

\begin{Shaded}
\begin{Highlighting}[]
\ExtensionTok{curl} \AttributeTok{{-}{-}fail} \AttributeTok{{-}{-}silent} \AttributeTok{{-}{-}show{-}error}\NormalTok{ http://localhost:18080/api/v1/health/live}
\end{Highlighting}
\end{Shaded}

Abrir \textbf{http://localhost:15173/}. Consultar las credenciales del
ambiente desarrollo en \texttt{CREDENCIALES\_LOCALES.txt} de forma
privada, sin compartir esa pantalla. El archivo contiene las claves
nuevas del paquete, no la contraseña habitual del tutor.

Para desarrollar se edita \texttt{fuente/backend} y
\texttt{fuente/frontend} dentro de la carpeta nueva. Esos son los
directorios montados por el Compose portable. Conservar la carpeta y su
ubicación mientras se utilice.

\clearpage

\section{Probar release y encender o apagar cada
ambiente}\label{probar-release-y-encender-o-apagar-cada-ambiente}

En un equipo de 8 GB, detener desarrollo antes de iniciar release. Los
dos ambientes tienen almacenamiento independiente; configurar uno no
configura automáticamente el otro.

\begin{Shaded}
\begin{Highlighting}[]
\ExtensionTok{docker}\NormalTok{ compose }\AttributeTok{{-}{-}env{-}file}\NormalTok{ config/desarrollo.env }\AttributeTok{{-}f}\NormalTok{ compose.desarrollo.yaml stop }\AttributeTok{{-}t}\NormalTok{ 60}
\end{Highlighting}
\end{Shaded}

\begin{Shaded}
\begin{Highlighting}[]
\ExtensionTok{docker}\NormalTok{ compose }\AttributeTok{{-}{-}env{-}file}\NormalTok{ config/release.env }\AttributeTok{{-}f}\NormalTok{ compose.release.yaml config }\AttributeTok{{-}{-}quiet}
\end{Highlighting}
\end{Shaded}

\begin{Shaded}
\begin{Highlighting}[]
\ExtensionTok{docker}\NormalTok{ compose }\AttributeTok{{-}{-}env{-}file}\NormalTok{ config/release.env }\AttributeTok{{-}f}\NormalTok{ compose.release.yaml up }\AttributeTok{{-}d} \AttributeTok{{-}{-}pull}\NormalTok{ never }\AttributeTok{{-}{-}wait} \AttributeTok{{-}{-}wait{-}timeout}\NormalTok{ 900}
\end{Highlighting}
\end{Shaded}

\begin{Shaded}
\begin{Highlighting}[]
\ExtensionTok{docker}\NormalTok{ compose }\AttributeTok{{-}{-}env{-}file}\NormalTok{ config/release.env }\AttributeTok{{-}f}\NormalTok{ compose.release.yaml ps }\AttributeTok{{-}a}
\end{Highlighting}
\end{Shaded}

\begin{Shaded}
\begin{Highlighting}[]
\ExtensionTok{docker}\NormalTok{ compose }\AttributeTok{{-}{-}env{-}file}\NormalTok{ config/release.env }\AttributeTok{{-}f}\NormalTok{ compose.release.yaml exec }\AttributeTok{{-}T}\NormalTok{ backend alembic current}
\end{Highlighting}
\end{Shaded}

Abrir \textbf{http://localhost:18088/} e ingresar con la credencial
\textbf{release} del paquete. Este modo sirve el frontend compilado y el
backend de entrega. No equivale por sí mismo a una promoción aprobada a
\texttt{main}: el manifiesto indicará la procedencia del código.

\begin{longtable}[]{@{}lll@{}}
\toprule\noalign{}
Uso & Desarrollo & Release \\
\midrule\noalign{}
\endhead
\bottomrule\noalign{}
\endlastfoot
Web local & \texttt{15173} & \texttt{18088} \\
API local & \texttt{18080} & \texttt{18081} \\
PostgreSQL publicado & \texttt{15432} & \texttt{15433} \\
SQL Server publicado & \texttt{11433} & \texttt{11434} \\
SQL desde el backend & \texttt{sqlserver:1433} &
\texttt{sqlserver:1433} \\
\end{longtable}

\subsection{Volver a desarrollo}\label{volver-a-desarrollo}

\begin{Shaded}
\begin{Highlighting}[]
\ExtensionTok{docker}\NormalTok{ compose }\AttributeTok{{-}{-}env{-}file}\NormalTok{ config/release.env }\AttributeTok{{-}f}\NormalTok{ compose.release.yaml stop }\AttributeTok{{-}t}\NormalTok{ 60}
\end{Highlighting}
\end{Shaded}

\begin{Shaded}
\begin{Highlighting}[]
\ExtensionTok{docker}\NormalTok{ compose }\AttributeTok{{-}{-}env{-}file}\NormalTok{ config/desarrollo.env }\AttributeTok{{-}f}\NormalTok{ compose.desarrollo.yaml up }\AttributeTok{{-}d} \AttributeTok{{-}{-}pull}\NormalTok{ never }\AttributeTok{{-}{-}wait} \AttributeTok{{-}{-}wait{-}timeout}\NormalTok{ 900}
\end{Highlighting}
\end{Shaded}

\subsection{Operación diaria}\label{operaciuxf3n-diaria}

Para apagar, utilizar el \texttt{stop\ -t\ 60} del ambiente en uso. Para
encender de nuevo, utilizar su \texttt{up} anterior. Estos comandos
conservan los volúmenes. No repetir extracción, \texttt{docker\ load},
limpieza de contenedores ni configuración inicial cada día.

No añadir \texttt{-\/-volumes} o \texttt{-v} para apagar. Los volúmenes
contienen configuración, usuarios, metadatos y datamarts; la clave de
cifrado debe conservarse junto con la base que protege.

\clearpage

\section{Configurar el sistema
limpio}\label{configurar-el-sistema-limpio}

Primero comprobar la URL del ambiente en uso. Repetir esta configuración
por separado en desarrollo y release si se utilizarán ambos.

\subsection{1. Registrar las dos conexiones de
demostración}\label{registrar-las-dos-conexiones-de-demostraciuxf3n}

En \textbf{Conexiones de datos}, crear una conexión para cada base:

\begin{longtable}[]{@{}
  >{\raggedright\arraybackslash}p{(\linewidth - 2\tabcolsep) * \real{0.5000}}
  >{\raggedright\arraybackslash}p{(\linewidth - 2\tabcolsep) * \real{0.5000}}@{}}
\toprule\noalign{}
\begin{minipage}[b]{\linewidth}\raggedright
Campo
\end{minipage} & \begin{minipage}[b]{\linewidth}\raggedright
Valor
\end{minipage} \\
\midrule\noalign{}
\endhead
\bottomrule\noalign{}
\endlastfoot
Motor & SQL Server \\
Servidor & \texttt{sqlserver} \\
Puerto interno & \texttt{1433} \\
Primera base & \texttt{AdventureWorks2022} \\
Segunda base & \texttt{WideWorldImporters} \\
Usuario lector & \texttt{bi\_reader} \\
Contraseña & La del lector SQL del ambiente activo en el archivo
privado \\
\end{longtable}

Para el certificado local, usar cifrado y confianza del certificado
cuando el formulario los ofrezca. No usar el puerto publicado del
computador como puerto interno Docker. No usar la clave de acceso web
como clave SQL ni trabajar con \texttt{sa} en la operación BI habitual.

En cada conexión: \textbf{Probar conexión → comprobar sólo lectura →
Habilitar → Leer metadatos}. Ambas pueden quedar habilitadas; el
selector de fuente determina el contexto de trabajo.

\subsection{2. Configurar OpenAI con consumo
controlado}\label{configurar-openai-con-consumo-controlado}

En \textbf{Configuración LLM}, crear una configuración con estos
valores:

\begin{itemize}
\tightlist
\item
  Proveedor: \textbf{OpenAI}.
\item
  URL: \texttt{https://api.openai.com/v1}.
\item
  Modelo inicial: \textbf{gpt-6-luna}.
\item
  Nivel de razonamiento: \textbf{Bajo}.
\end{itemize}

Es el preset económico vigente; su acceso y disponibilidad se
comprobarán con la cuenta de API utilizada. El campo modelo sigue siendo
editable.

Guardar el registro; después usar \textbf{Registrar credencial} e
ingresar la API key directamente en el campo privado. No escribirla en
Parámetros, código, documentos, capturas ni chats. Usar \textbf{Probar
conexión} y activar la configuración cuando la prueba resulte
satisfactoria.

El paquete limpio no lleva la clave ni el crédito del tutor. Las
llamadas reales requieren Internet y consumen la cuota de la cuenta
configurada. Vigilar el saldo y evitar repetir generaciones sólo para
obtener capturas. Una prueba de conexión no demuestra por sí sola la
generación correcta de una propuesta BI.

\subsection{3. Revisar duración de
sesión}\label{revisar-duraciuxf3n-de-sesiuxf3n}

Abrir \textbf{Parámetros}, localizar el grupo \textbf{Seguridad} y el
campo \textbf{Duración de la sesión (minutos)}.

Su identificador es \texttt{AUTH\_SESSION\_MINUTES}. El valor inicial es
\textbf{60} y el rango permitido es \textbf{15--240} minutos. Guardar el
valor elegido y volver a iniciar sesión. El cambio afecta los tokens
nuevos; no alarga uno ya emitido.

\subsection{4. Preparar el catálogo de cada
fuente}\label{preparar-el-catuxe1logo-de-cada-fuente}

En \textbf{Explorador de esquema}, revisar la cabecera y la instantánea
de la fuente elegida. Después, en \textbf{Catálogo analítico}, revisar
preguntas y periodicidades de esa misma conexión y guardar. Repetir con
la otra base. Las preguntas personalizadas no se copian entre bases por
coincidencia de nombre.

\clearpage

\section{Recorrer BI con ambas bases y registrar la
prueba}\label{recorrer-bi-con-ambas-bases-y-registrar-la-prueba}

La instalación limpia genera sus propios identificadores. No buscar
números de propuestas o ejecuciones que aparezcan en capturas del tutor.

\subsection{Primera fuente}\label{primera-fuente}

\begin{enumerate}
\def\labelenumi{\arabic{enumi}.}
\tightlist
\item
  Elegir AdventureWorks en la cabecera; confirmar base e instantánea.
\item
  Revisar tablas, claves y relaciones en el Explorador. El evento
  comercial de los ejemplos de pedidos debe conservar esa
  interpretación.
\item
  Revisar el catálogo y seleccionar preguntas pertinentes.
\item
  En el Asistente, describir una necesidad acotada, por ejemplo comparar
  mensualmente importes, unidades y pedidos por producto y territorio,
  con promedio por pedido.
\item
  Autorizar el destino LLM mostrado y validar la viabilidad de la
  necesidad. Resolver ambigüedades; el esquema no demuestra por sí solo
  significado comercial, moneda ni calidad de filas.
\item
  Generar la propuesta y recorrer Necesidad, Conceptos, Propuesta,
  Personalización y Revisión. Comprobar granularidad, medidas, KPI y
  relaciones. Aprobar sólo el resultado compatible.
\item
  En Datamart de ventas: elegir propuesta, revisar los indicadores,
  confirmar y materializar. La aprobación anterior no ejecutó el ETL.
\item
  Revisar conciliación, cantidades, importes, identidad descriptiva y
  posibles decisiones semánticas antes de interpretar los resultados.
\item
  En Analítica, comprobar fuente, ejecución y propuesta; aplicar año y
  territorio; revisar el denominador de cada promedio. Exportar PDF y
  Excel y contrastarlos con la selección visible.
\end{enumerate}

\subsection{Segunda fuente}\label{segunda-fuente}

Repetir el recorrido con WideWorldImporters. Para un caso de
facturación, solicitar expresamente ventas facturadas, líneas de
factura, comprador y ubicación relacionada pertinente. No convertir
pedidos en facturas ni confundir comprador con destinatario de
facturación o ciudad postal con ciudad de entrega.

\subsection{Comprobar continuidad}\label{comprobar-continuidad}

\begin{itemize}
\tightlist
\item
  Recargar un módulo: debe conservarse ese módulo autorizado.
\item
  Entrar desde Inicio o Revisar configuración: \textbf{Atrás} debe
  volver al recorrido que lo abrió, conservando el paso. Una entrada
  directa desde la barra no necesita ese botón contextual.
\item
  En Analítica, usar Consultar respaldo y volver: conservar ejecución,
  indicador, vista y filtros confirmados.
\item
  Alternar las dos fuentes y regresar: cada conexión conserva su propia
  selección dentro de la pestaña, sin traer filtros, datos o chat de la
  otra.
\item
  Si una pregunta del borrador fue retirada del catálogo, el objetivo se
  conserva y se pide revisar la selección y validar otra vez. Las
  propuestas históricas conservan su evidencia.
\item
  Un KPI añadido por una regla desde un requisito validado debe explicar
  que procede de la plataforma; no atribuir esa adición al proveedor IA.
\end{itemize}

Registrar fecha, equipo, ambiente, corte del manifiesto, fuente, IDs
propios, resultado y captura sin secretos. Los ensayos nuevos de un
modelo deben anotarse como tales: los resultados observados con otro
modelo no los certifican automáticamente.

\clearpage

\section{Resolver los incidentes que ya
ocurrieron}\label{resolver-los-incidentes-que-ya-ocurrieron}

\subsection{El directorio anterior fue borrado con
Supr}\label{el-directorio-anterior-fue-borrado-con-supr}

Los contenedores y volúmenes pueden seguir existiendo. Utilizar los
pasos de identificación y retirada por etiquetas. No extraer una versión
nueva dentro de la carpeta antigua desaparecida para intentar
recuperarla ni usar un Compose nuevo con un proyecto viejo.

\subsection{Error: puerto 11433 ya
ocupado}\label{error-puerto-11433-ya-ocupado}

\begin{Shaded}
\begin{Highlighting}[]
\ExtensionTok{docker}\NormalTok{ ps }\AttributeTok{{-}{-}format} \StringTok{\textquotesingle{}table \{\{.Names\}\}\textbackslash{}t\{\{.Image\}\}\textbackslash{}t\{\{.Ports\}\}\textquotesingle{}}
\end{Highlighting}
\end{Shaded}

Buscar el contenedor que publica ese puerto. En la experiencia anterior
era \texttt{bi-limpio-dev-sqlserver-1}. Si pertenece al BI viejo,
aplicar su retirada selectiva. Si es otro proyecto, detener esta
sustitución y revisar: no eliminarlo. Después repetir el \texttt{up} del
ambiente nuevo.

\subsection{Error: migrate terminó con código
1}\label{error-migrate-terminuxf3-con-cuxf3digo-1}

\begin{Shaded}
\begin{Highlighting}[]
\ExtensionTok{docker}\NormalTok{ compose }\AttributeTok{{-}{-}env{-}file}\NormalTok{ config/desarrollo.env }\AttributeTok{{-}f}\NormalTok{ compose.desarrollo.yaml logs }\AttributeTok{{-}{-}tail}\NormalTok{ 120 migrate}
\end{Highlighting}
\end{Shaded}

\begin{Shaded}
\begin{Highlighting}[]
\ExtensionTok{docker}\NormalTok{ compose }\AttributeTok{{-}{-}env{-}file}\NormalTok{ config/desarrollo.env }\AttributeTok{{-}f}\NormalTok{ compose.desarrollo.yaml logs }\AttributeTok{{-}{-}tail}\NormalTok{ 80 postgres seed backend}
\end{Highlighting}
\end{Shaded}

La causa se determina leyendo ese mensaje, no sólo el estado rojo.
Comprobar que se usa la carpeta, configuración y proyecto nuevos; que la
migración actual está incluida y que no se reutilizó un volumen antiguo
con otra configuración. No cambiar contraseñas por aproximación, no
editar \texttt{alembic\_version} y no borrar volúmenes como primer
intento.

\subsection{SQL Server o la espera de 900 segundos no
termina}\label{sql-server-o-la-espera-de-900-segundos-no-termina}

\begin{Shaded}
\begin{Highlighting}[]
\ExtensionTok{docker}\NormalTok{ compose }\AttributeTok{{-}{-}env{-}file}\NormalTok{ config/desarrollo.env }\AttributeTok{{-}f}\NormalTok{ compose.desarrollo.yaml logs }\AttributeTok{{-}{-}tail}\NormalTok{ 100 sqlserver}
\end{Highlighting}
\end{Shaded}

\begin{Shaded}
\begin{Highlighting}[]
\FunctionTok{free} \AttributeTok{{-}h}
\end{Highlighting}
\end{Shaded}

\begin{Shaded}
\begin{Highlighting}[]
\FunctionTok{df} \AttributeTok{{-}h} \StringTok{"}\VariableTok{$HOME}\StringTok{"}
\end{Highlighting}
\end{Shaded}

Comprobar que sólo está activo un ambiente, que existe espacio y que la
restauración sigue avanzando. Un segundo \texttt{up} puede retomar la
preparación idempotente; no obliga a eliminar datos. Para release, usar
sus dos archivos en las mismas consultas.

\subsection{La propuesta o el proveedor
falla}\label{la-propuesta-o-el-proveedor-falla}

Leer el mensaje antes de reintentar. Comprobar catálogo visible, fuente,
autorización de metadatos, modelo, credencial, prueba y saldo. Un
rechazo del contrato no se corrige aceptando cualquier JSON ni
inventando columnas. Conservar la necesidad y la evidencia, resolver la
causa y generar una versión nueva cuando corresponda. No repetir
llamadas costosas de forma automática.

Los registros se revisan localmente; ocultar datos privados antes de
compartirlos. No imprimir \texttt{docker\ compose\ config} completo,
\texttt{.env}, credenciales o inspecciones Docker completas en capturas.

\clearpage

\section{Qué se conserva y cuándo está
listo}\label{quuxe9-se-conserva-y-cuuxe1ndo-estuxe1-listo}

\subsection{Código de desarrollo y código de
release}\label{cuxf3digo-de-desarrollo-y-cuxf3digo-de-release}

El portable utiliza \textbf{una fuente de autoría} en \texttt{fuente/}.
Desarrollo la monta y permite trabajar en ella. Release utiliza imágenes
compiladas del corte declarado; editar \texttt{fuente/} no cambia esas
imágenes. Una nueva entrega release exige reconstruir y comprobar sus
imágenes desde el código elegido.

Los clones viejos
\texttt{\textasciitilde{}/Proyectos/prototipo-bi-ia-dev} y
\texttt{\textasciitilde{}/Proyectos/prototipo-bi-ia-release} no
actualizan este portable ni son la fuente que monta su Compose. No
iniciar un tercer conjunto desde \texttt{fuente/compose.yaml} o mediante
\texttt{make\ up} durante esta instalación. Usar los Compose situados en
la raíz del paquete.

\subsection{Bases, volúmenes y
cifrado}\label{bases-voluxfamenes-y-cifrado}

\texttt{docker\ load} importa \textbf{imágenes}, no los volúmenes del
tutor. PostgreSQL obtiene su estructura por migraciones y semilla. SQL
Server obtiene las muestras públicas desde el respaldo incluido en la
imagen offline. Los cambios realizados después del inicio viven en los
volúmenes del ambiente receptor.

La clave raíz de cifrado se crea en su volumen local
\texttt{secret\_key\_data}. Protege las credenciales guardadas por la
aplicación. La base PostgreSQL y esa clave deben mantenerse juntas para
recuperar secretos existentes. No importar sólo PostgreSQL desde otra
instalación ni copiar una clave del tutor sobre la recién creada.

\subsection{Lista de aceptación del
receptor}\label{lista-de-aceptaciuxf3n-del-receptor}

\begin{itemize}
\tightlist
\item[$\square$]
  Sumas externa e internas correctas; imágenes coinciden con el
  manifiesto nuevo.
\item[$\square$]
  Ningún contenedor o volumen ajeno al BI antiguo fue retirado.
\item[$\square$]
  Desarrollo y release se probaron por separado; cuatro servicios
  saludables y \texttt{migrate}/\texttt{seed} con salida cero.
\item[$\square$]
  Migración 20261125\_20 o posterior declarada presente; sesión
  parametrizable comprobada.
\item[$\square$]
  Administrador accede y utiliza credenciales propias del ambiente.
\item[$\square$]
  Ambas muestras conectan en sólo lectura y muestran sus metadatos.
\item[$\square$]
  Catálogo, necesidad, propuesta, ETL y analítica conservan la fuente
  correcta.
\item[$\square$]
  Conciliación y exportaciones revisadas; filtros y ejecución
  conservados al volver a una fuente.
\item[$\square$]
  IA probada sólo con credencial autorizada y saldo comprobado;
  limitaciones anotadas.
\item[$\square$]
  Después de \texttt{stop} y \texttt{up}, se conserva el trabajo del
  receptor.
\item[$\square$]
  Carpeta nueva conservada; carpetas antiguas retiradas sólo al concluir
  la comprobación.
\end{itemize}

En Windows, la instalación nueva se realiza desde Ubuntu en WSL 2 con
Docker compatible y los archivos dentro del HOME Linux. Las órdenes
destructivas protegidas para \texttt{ls-va@Lenovo} \textbf{no se adaptan
automáticamente a ese equipo}. Revisar sus requisitos y usuario propios
antes de una sustitución; la otra programadora indicó una instalación
inicial, no una limpieza de proyectos ajenos.

\subsection{Base documental y alcance de esta
edición}\label{base-documental-y-alcance-de-esta-ediciuxf3n}

Esta guía toma como antecedente la instalación directa y la guía
portable del 17-11-2026. Sus comandos de importación, inicio, parada y
consulta se conservan, ahora con nombres leídos del paquete definitivo y
retirada individual por etiquetas aunque el directorio antiguo falte. Se
incorporan las correcciones operativas y la configuración económica
OpenAI del 25-11-2026.

La elaboración de este documento no construyó el paquete nuevo ni
ejecutó instalaciones o retiradas en el equipo receptor. El manifiesto y
el informe de aceptación del nuevo portable deberán registrar sus
propios resultados antes de distribuirlo.

\end{document}
